\section{A trace-coordinate check of the WP factor}
\label{app:fricke}

A separate two-dimensional calculation is useful for checking the factor
$2$ in the local densities. Let $A,B,C,D\in\mathrm{SL}_2(\RR)$ satisfy
$ABCD=I$, and set
\[
 a=\tr A,\ b=\tr B,\ c=\tr C,\ d=\tr D,
 \qquad x=\tr(AB),\ y=\tr(AC),\ z=\tr(BC).
\]
The matrices are holonomies around the four boundary loops in
their chosen order; $ABCD=I$ expresses the surface relation.
The trace of a matrix is the sum of its diagonal entries.
Peripheral traces $a,b,c,d$ are held fixed, whereas $x,y,z$
describe internal curves and vary with the marked structure.
One relation among those three coordinates leaves a surface.

With the four peripheral traces fixed, the character surface satisfies
the Fricke relation
\begin{align}
 \Phi(x,y,z)={}&x^2+y^2+z^2+xyz-(ab+cd)x-(ac+bd)y\notag\\
 &-(ad+bc)z+a^2+b^2+c^2+d^2+abcd-4=0.
 \label{eq:fricke}
\end{align}
This is the standard four-holed-sphere trace surface; see~\cite{Goldman}.

On a chart where $x=2\cosh(\ell/2)>2$, put $s=\sinh(\ell/2)$ and
$X=AB$. Twisting across the $AB$ cuff fixes $A,B$ and conjugates $C,D$ by
\[
 \exp\left(\frac{\tau}{2}H_X\right),
 \qquad H_X=\frac{X-(x/2)I}{s}.
\]
The eigenvalues of $H_X$ are $1,-1$, so this is a displacement of length
$\tau$ along the cuff axis. Differentiating $y=\tr(AC)$ gives
\[
 \frac{\partial y}{\partial\tau}
 =\frac1{2s}\tr\bigl(A[X,C]\bigr)
 =-\frac{\Phi_z}{2s}.
\]
For the last equality, Cayley--Hamilton gives
$\tr(A^2BC)=ad-z$, while the two-matrix trace identity gives
$\tr(ACAB)=xy-bc+z$. Their difference is
$ad+bc-xy-2z=-\Phi_z$.

Here $[X,C]=XC-CX$ is a matrix commutator and
$\Phi_z=\partial\Phi/\partial z=2z+xy-ad-bc$.
The calculation is a change of variables from length and twist
$(\ell,\tau)$ to trace coordinates $(x,y)$.
Its Jacobian is $s(-\Phi_z/(2s))=-\Phi_z/2$:
the factors of $\sinh(\ell/2)$ cancel, leaving the desired two.

Since $\dd x=s\,\dd\ell$ and $\omega_{\WP}=\dd\ell\wedge\dd\tau$,
\begin{equation}
 \boxed{\dV=\frac{2|\dd x\wedge\dd y|}{|\Phi_z|}.}
 \label{eq:frickedensity}
\end{equation}
As a local illustration, at a point with $|\Phi_z|=4$ a tiny
trace-coordinate rectangle of side lengths $\epsilon,\eta$
has WP area $\epsilon\eta/2+o(\epsilon\eta)$ as both sides
shrink. This uses the density at an already geometric smooth
point; it does not construct such a point from arbitrary traces.

Equivalently, on this twist chart,
$\dd x\wedge\dd y=-\Phi_z\,\dd\ell\wedge\dd\tau/2$.
Reversing the chosen twist direction reverses the signed form but leaves
the displayed positive density unchanged. The hyperbolic application
requires the smooth geometric component of the relative character
surface and $\Phi_z\ne0$ for $(x,y)$ to be coordinates. If $\Phi_z=0$
but $\nabla\Phi\ne0$, another pair can be chosen; if all three partial
derivatives vanish, the character surface is singular and this inversion
argument does not supply a chart. In particular, arbitrary real solutions
of $\Phi=0$ are not automatically holonomies of four-holed spheres.
The formula is local on the marked character surface and does not
describe a global mapping-class fundamental domain.

